\subsection{LLM Agents and Reasoning Frameworks}

Large Language Models (LLMs) have recently demonstrated strong capabilities beyond traditional natural language processing tasks. Modern LLMs are capable of performing reasoning, planning, and sequential decision-making when provided with structured prompts. This has led to the emergence of \textbf{LLM agents}, where language models are placed in an interaction loop with an environment and generate actions based on observations.

Formally, an LLM agent operates within a closed-loop interaction cycle consisting of three stages:

% \begin{center}
% Observation $\rightarrow$ Reasoning $\rightarrow$ Action
% \end{center}

\begin{tikzpicture}[
    node distance=2.5cm,
    box/.style={rectangle, draw, rounded corners, minimum width=2.5cm, minimum height=1cm, align=center},
    arrow/.style={-{Stealth[length=3mm]}, thick}
]
    \node[box] (obs)    {Observation};
    \node[box, right=of obs] (reason) {Reasoning};
    \node[box, right=of reason] (act)    {Action};

    % Forward arrows
    \draw[arrow] (obs) -- (reason);
    \draw[arrow] (reason) -- (act);

    % Loop-back arrow from Action to Observation
    \draw[arrow] (act.south) -- ++(0,-1) -| (obs.south);
\end{tikzpicture}

At each timestep, the environment produces an observation describing the current state. This observation is converted into a textual representation and provided to the language model as part of a prompt. The LLM processes the prompt and generates a response that contains the next action to perform. The selected action is then executed in the environment, which updates the state and produces the next observation.

Recent research has shown that \textbf{prompting strategies} can significantly influence the reasoning behavior of LLM agents. Instead of directly predicting an answer or action, models can be encouraged to generate intermediate reasoning steps. Techniques such as \textbf{Chain-of-Thought (CoT)} prompting demonstrate that explicitly asking the model to reason step-by-step can improve performance on complex tasks. Subsequent work has extended this idea to more sophisticated reasoning frameworks such as \textbf{ReAct}, \textbf{Tree-of-Thought}, and \textbf{Buffer-of-Thought}, which incorporate structured reasoning, branching exploration, and reusable reasoning templates.

These reasoning frameworks are particularly relevant for \textbf{sequential decision-making tasks}, where agents must select a sequence of actions while considering environmental constraints. Navigation tasks in grid environments provide a controlled setting to evaluate these reasoning strategies. In such environments, the agent must reach a goal location while avoiding hazards such as walls or lava tiles. The agent must interpret spatial information, evaluate potential moves, and plan safe paths toward the goal.

In this work, we evaluate different reasoning strategies for LLM agents in navigation tasks using the MiniGrid environment. Specifically, we compare reactive decision-making with reasoning-based approaches including Chain-of-Thought, ReAct, Tree-of-Thought, and Buffer-of-Thought. By evaluating these methods across environments of increasing difficulty, we aim to understand the trade-offs between reasoning complexity, computational cost, and navigation performance.